% Week 8 Session A — Testing & Validation
% Upload: Week8_SessionA_Report.tex, week8_session_a_terminal.png, week8_session_a_files/
%
\documentclass[11pt,a4paper]{article}

\usepackage[margin=2.5cm]{geometry}
\usepackage{graphicx}
\newcommand{\termfig}[1]{%
  \includegraphics[width=\linewidth,height=0.55\textheight,keepaspectratio]{#1}%
}
\usepackage{booktabs}
\usepackage{xurl}
\usepackage{hyperref}
\usepackage{parskip}
\usepackage{enumitem}
\usepackage{xcolor}
\usepackage{fvextra}
\usepackage[most]{tcolorbox}

\newcommand{\studentname}{Mahmudul Hasan}
\newcommand{\studentid}{4125999049}
\newcommand{\reportdate}{May 22, 2026}
\newcommand{\capstonerepo}{https://github.com/mahmud456alhasan-debug/smart-water-capstone}

\makeatletter
\g@addto@macro\UrlBreaks{\do\ }
\makeatother

\newcommand{\LabCodeRoot}{week8_session_a_files}

\newcommand{\filebox}[2]{%
\begin{tcolorbox}[
  enhanced,
  colback=gray!6,
  colframe=black!55,
  title={#1},
  fonttitle=\bfseries\ttfamily\footnotesize,
  arc=1.5mm,
  boxrule=0.7pt,
  breakable,
  width=\linewidth,
  before skip=0.8em,
  after skip=0.8em,
]
\VerbatimInput[fontsize=\footnotesize,breaklines=true,breakanywhere=true]{\LabCodeRoot/#2}
\end{tcolorbox}%
}

\title{Lab Report: Week 8 Session A\\Testing \& Validation}
\author{\studentname \quad (Student ID: \studentid)}
\date{\reportdate}

\begin{document}
\maketitle

\begin{abstract}
This report documents Week~8 Session~A on the Smart Water capstone: \textbf{29 pytest cases}, \textbf{96\%} statement coverage on \texttt{src/}, a physical validation module (\texttt{src/validation.py}), and integration workflows (Swiss Cheese Model). A hallucination in an edge-case assumption ($P{=}5$~mm always below $I_a$ for high CN) was corrected after tests failed. Evidence: Figure~\ref{fig:terminal} (\texttt{pytest --cov}). Repository: \url{\capstonerepo}.
\end{abstract}

\section{Introduction}
After Week~7 implementation, the capstone had only two unit tests (27\% coverage). Session~A expands testing across runoff, weather, reservoir, flood, validation, and integration layers per the lab guide.

\section{Objectives}
\begin{itemize}[noitemsep]
  \item Comprehensive unit tests (Exercise~1).
  \item Swiss Cheese verification layers (Exercise~2).
  \item Find and fix AI/human test assumptions (Exercise~3).
  \item Document in \texttt{prompt\_log.md} (Exercise~4).
\end{itemize}

\section{Exercise 1: Comprehensive tests}

\subsection{Test inventory}
\begin{table}[htbp]
  \centering
  \caption{Test modules and focus.}
  \label{tab:tests}
  \small
  \begin{tabular}{@{}ll@{}}
    \toprule
    File & Coverage \\
    \midrule
    \texttt{test\_runoff.py} & Normal storm, $P{=}0$, $Q \le P$, invalid CN/$P$ \\
    \texttt{test\_weather.py} & CSV load, missing columns, GREEN/AMBER/RED \\
    \texttt{test\_flood.py} & Mask, depth, DEM bounds, monotonicity \\
    \texttt{test\_reservoir.py} & Revenue $\approx$\$708{,}849, storage PASS \\
    \texttt{test\_validation.py} & Physical check failures \\
    \texttt{test\_integration.py} & Rainfall$\rightarrow$alert, full workflows \\
    \bottomrule
  \end{tabular}
\end{table}

\subsection{Coverage results}
Before Session~8: \textbf{2 passed}, \textbf{27\%} coverage. After: \textbf{29 passed in 0.63s}, \textbf{96\%} on \texttt{src/} (201 statements, 8 missed---mostly DEM auto-generate branch and rare optimizer error paths).

\begin{figure}[htbp]
  \centering
  \termfig{week8_session_a_terminal.png}
  \caption{Terminal: \texttt{pytest -q --cov=src --cov-report=term-missing} --- 29 passed, 96\% total coverage.}
  \label{fig:terminal}
\end{figure}

\section{Exercise 2: Swiss Cheese Model}

\begin{table}[htbp]
  \centering
  \caption{Verification layers applied.}
  \label{tab:cheese}
  \begin{tabular}{@{}cl@{}}
    \toprule
    Layer & Action \\
    \midrule
    1 & Manual code review (units, Week~6 reservoir logic) \\
    2 & 29 unit tests, all passing \\
    3 & \texttt{src/validation.py}: $Q \le P$, storage $[50,200]$ MCM, flood monotonicity \\
    4 & \texttt{tests/test\_integration.py}: CSV$\rightarrow$alert, DEM$\rightarrow$stats, reservoir PASS \\
    \bottomrule
  \end{tabular}
\end{table}

\section{Exercise 3: Hallucination / assumption fix}

\textbf{Issue:} Test assumed 5~mm rainfall always yields zero runoff for CN${=}95$. High CN lowers $I_a$, so computed $Q \approx 0.34$~mm.

\textbf{Fix:} Use $P{=}2$~mm with CN${=}80$ for a true below-$I_a$ case.

\textbf{Lesson:} Physical edge cases must use formula-derived $I_a$, not intuition.

\section{Exercise 4: Repository}
Updated tests and \texttt{src/validation.py} committed to \url{\capstonerepo}. \texttt{prompt\_log.md} includes Session~8 entries (Appendix~\ref{app:promptlog}).

\section{Discussion}
\begin{enumerate}[noitemsep]
  \item Coverage jumped from 27\% to 96\% by testing weather and reservoir modules, not only runoff stubs.
  \item Integration tests catch tab-level workflows without running Streamlit headless.
  \item One failing test is valuable---it exposed a weak edge-case assumption, not production code failure.
\end{enumerate}

\section{Conclusion}
Week~8 Session~A met the lab guide: broad pytest suite, $>$80\% coverage, Swiss Cheese layers, documented fix, and terminal evidence. Capstone remains deployable via \texttt{streamlit run app/main.py} with validated hydrologic constraints.

\appendix

\section{Physical validation: \texttt{validation.py}}
\label{app:validation}
\filebox{\texttt{src/validation.py} (full file)}{validation.py}

\section{Submitted prompt log (excerpt)}
\label{app:promptlog}
\filebox{\texttt{prompt\_log.md} (Session 8 section)}{prompt_log_session8.md}

\end{document}
